
\chapter{Méthodologie}

\maketitle

\section{Introduction}


\subsection{problèmes du nombre de parcelles}

Comme nous l’avons vu, le fichier Excel Barbéro contient 17803 parcelles.

Le numéro maximal attribué à la dernière parcelle répertoriée dans les documents de la Mappe de Beaufort aux ADS est : 17691.

Enfin le nombre de numéros de parcelles géoréférencées extraites des cartes Barbéro est seulement de 11787.

Une étape préalable consiste donc à comprendre ces différences.

a)	Nous avons constaté dans le fichier Beaufort-FICHIER-ATLAS-OK.xls qu’il manquait 136 numéros dans la séquence continue de 1 à 17691 ; en consultant les registres aux ADS, et en effectuant un sondage, il s’avère que ces numéros « manquants » y sont rayés et il est fait mention que ces parcelles ont été transférées  à la commune de « La Chapelle » à un moment donné durant l’élaboration de la Mappe ; « La vallée de Treicol est un secteur âprement disputé par Aime et les Chapelles, contre Saint-Maxime. Elle fait l’objet d’un procès pluriséculaire, puisqu’il débute en 1599 par le refus de Me Jean-François Dunand de payer la taille de sa montagne de Presset à Beaufort. L’affaire s’envenime, car de nombreux particuliers possèdent des alpages à Treicol. Pour être limitrophe de la vallée et augmenter ses prétentions, la communauté des Chapelles va jusqu’à acheter en 1732 une longue et étroite bande de rochers depuis la Pierre-Menta jusqu’à son territoire. Portée devant toutes les instances possibles, l’affaire n’est définitivement tranchée qu’en 1907, par un arrêt du Conseil d’État qui attribue la vallée à Beaufort.» 	\footcite{HViallet1993}.

La Chapelle est devenue la commune actuelle de la cote d’Aime (fusionnée en 2016 dans la Plagne Tarentaise). 

Carte9 : la zone « inconnu » correspond à l’information du « mas » des parcelles, ces parcelles n’étant pas référencées dans le fichier Excel Barbéro. 

Nous avons ajouté au fichier les 136 parcelles manquantes, ces lignes ne comportent aucun attribut descriptif.

b)	Autre particularité du fichier Excel Barbéro est qu’il contient :

-	Une parcelle non numérotée, mais avec la lettre « A » dans le champ n° de parcelle.

-	6 parcelles différentes, mais dont les numéros sont en doublon.


Récapitulatif :

Analyse des numéros 
du tb excel Barbéro	Tableau
Barbéro	Nb Entiers
+ Parcelles
manquantes	Hypothèses
Parcelles avec nb entier entre 1 et 17691	17552	17552	
Parcelles avec le même n° (nombre entier)	2	1	le doublon est une décimale omise
Parcelles manquantes (Transférées)		136	
Parcelle numérotée "A'	1	0	le n° omis est un nombre décimal
Parcelles avec un n° + une lettre	2	2	
Parcelles avec n° + lettre (n° déjà existant)	10		
Parcelles avec une décimale (1/2 ; 1/3,…)	232		
Parcelles avec le même n° nb + décimale)	4		
Total	17803	17691	

Conclusion : si nous totalisons uniquement les numéros avec un nombre entier, nous retombons sur le dernier numéro de parcelle (17691) des registres de la Mappe de Beaufort.


\subsection{Import de la Mappe dans QGIS}

La première étape \footnote{projet QGIS : Assemblage précis.qgz} a consisté à superposer à plat les différents éléments disponibles Cette approche évite à ce stade le  risque de déformation du fait d'un géoréférencement d'une carte ancienne sur une carte moderne IGN) : 

- calage les unes par rapport aux autres, les quatre images de la vue globale de la Mappe Sarde

- superposition par-dessus (avec l'outil Georeférencer de QGIS) les 45 photos détaillées de la Mappe téléchargées des ADS. 



\cartefigure
{figures/1 MS 45 photos détails non GeoRef.jpg}
{fig:MS45nonGeoRef}
{45 photos assemblées de la Mappe Sarde de Beaufort}

\cartefigure
{figures/1 Assemblage des 45 photos détail Mappe Sarde Beaufort.jpg}
{fig:MS45Assemblage}
{assemblage des 45 photos de la Mappe Sarde de Beaufort}

Remarques :\\
- une photo est manquante sur le site des ADS (\texttt{361\_b\_16}). Afin de la constituer, nous avons découpé un morceau de  la photo globale associée . Bien évidemment la qualité est beaucoup plus faible.

- Le collage des 45 photos a fait apparaître quelques non-recouvrements du fait du photographe des ADS

- la Mappe Sarde de Beaufort est parfois très abîmée en particulier la jonction entre les deux toiles peintes est impossible du fait des manques.

\cartefigure
{figures/1 MS 45 photos détails Anomalies.jpg}
{fig:MS45_Anomalies}
{photos 361_013 et 014 non jointives avec 019 et 020}






\subsection{Retraitement des cartes Barbéro}

Le CD-Rom fournit 55 images-cartes comme celle-ci :

\cartefigure[0.5\textwidth]
{figures/1 exemple carte CD Barbéro.JPG}
{fig:BarberoCarte01}
{exemple Carte 01 CD Barbéro} 

Telle quel, ces cartes-images, ne sont pas plus exploitable que les photos de la Mappe déjà importées.
En revanche, elles contiennent deux éléments seulement : des lignes et des chiffres.

Des outils permettent des retraitements suffisamment automatiques pour rendre possible une alternative au travail manuel sur les 55 images et les 17691 parcelles. Le détail des outils et des étapes est décrit en annexe (cf Annexe3), en voici les tâches principales :

-	Détourage des 55 photos pour ne garder que les lignes et les numéros

-	Séparation des figures : 55 images avec seulement les contours des parcelles et 55 images avec les numéros des parcelles seulement .

\cartefigure[0.5\textwidth]
{figures/1 Beaufort secteur 01 original.jpg}
{fig:BarberoCarteNumOnly}
{Détourage et conservation des numéros exemple Carte 01}


\cartefigure[0.5\textwidth]
{figures/1 Beaufort secteur 01 reds.jpg}
{fig:BarberoCarteRed}
{Détourage et conservation des contours des parcelles exemple Carte 01}



-	Décodage des chiffres, les outils d’OCR que nous avons essayés ne donnant pas de résultat, nous avons dû coder en python une reconnaissance de formes au niveau pixel ; chaque un des nombres (0 à 9) ayant plus ou moins une signature en pixel unique.

-	Nous avons gardé la position des ensembles de pixels composant chaque chiffre. Puis nous avons regroupé les chiffres mitoyens pour aboutir a) à un nombre = au numéro de la parcelle et b) une référence (X,Y) du pixel au centre du groupe de chiffre composant un nombre.

-	Un retraitement manuel sur une centaine de nombres non reconnus a été nécessaire : nombre tronqué en bordure d’image, nombre mélangé à des éléments superposés comme la boussole ou la barre d’échelle ou encore chiffre mal pixélisé non reconnu dans le traitement précédent.

Au final de cette étape nous obtenons :

-	55 images des contours de parcelles 

- 	55 images des numéros (images inutilisées sauf pour vérification)

-	55 tables avec les nombres reconnus et leur position pixel (X,Y) sur chaque image.

Les 3 ensembles sont exactement superposables

Nous avons rassemblé les 55 tables en un seul fichier \texttt{XY\_No\_Parc.csv} (fichier ; X ; Y ; No\_Parc) et nous avons éliminé les doublons (parcelle se trouvant affichée sur deux images de secteur mitoyen). Au final nous avons 11.744 numéros distincts reconnus. Ce nombre représente les 2/3 (66\%) des 17.803 parcelles référencées dans la table Barbéro.

C'est assez compréhensible : la table Barbéro (fichier Excel) est issue des registres associés à la Mappe alors que les numéros sur les cartes sont issus de la lecture de la carte elle-même où les numéros ne sont pas toujours lisibles.

Nous verrons comment localiser approximativement les parcelles du tiers restant avec une méthode fondée sur la localisation du Mas de chaque parcelle. 




\section{Attribution d’une altitude à chaque parcelle}

\subsection{Géoréférencement des cartes Barbéro}

Le CD-Rom contient une carte d’assemblage des 55 secteurs :


\cartefigure 
{figures/1 Tableau d'assemblage des secteurs.JPG}
{fig:BarberoAssemblage}
{Carte d'assemblages de cartes fournie dans le CD Barbéro}


En utilisant la fonction géoréférencement de QGIS, nous avons superposé les 55 images des contours des parcelles sur l’assemblage détaillé de la Mappe. Cette opération était relativement simple (mais longue !) puisque nous pouvons assez facilement faire correspondre des points de ce géoréférencement en retrouvant les lignes des parcelles des cartes Barbéro avec celles correspondantes de la Mappe.

Malheureusement dans un ou deux cas l’assemblage des cartes Barbéro comportait des lacunes (certaines cartes ne se superposaient pas). C’est dommage, mais l’impact est resté très marginal.

Ensuite, nous avons réutilisé les mêmes points de géoréférencement de l’étape précédente, pour géoréférencer les cartes des numéros de parcelles, afin que chaque numéro se retrouve au bon endroit.

À ce stade, nous avons 3 couches superposées avec :\\
- les photos détaillées de la Mappe\\
- les contours des parcelles\\
- les numéros (en mode image, inexploitables mais utile pour vérifications ultérieures)


La dernière étape a consisté à transformer la table des numéros de parcelles en une couche vectorielle avec une table d’attributs obtenant ainsi la géolocalisation des numéros de parcelles.

un programme python qui : \\
- prend chaque la table "i"\\
- la charge en layer vecteur (add delimited text)\\
- l'exporte en fichier gpkg (format de couche vectorielle de QGIS)\\
-charge ce fichier dans le GeoRefrencer de QGIS et l'associe au fichier des points utilisé pour géolocaliser les contour des parcelles\\
- lance le géoréférencement

Au final nous avons 55 fichiers gpkg qui se superposent aux trois couches raster précédentes.

Pour chacune des couches je fusionne tous les éléments en un seul :\\
- les 45 images raster de la Mappe en un raster virtuel : \texttt{Virtual photo detaillées.vrt} \\
- les 55 images des contours des parcelles en un raster virtuel \texttt{Virtual_Red.vrt}\\
- les 55 fichiers gpkg des numéros de parcelles en un seul gpqk  \texttt{text}\\

Remarque : il y a dans le fichier des numéros de parcelles des numéros en double (ou triple et même quadruple. Ce sont sans doute des numéros qui apparaissent sur plusieurs secteurs Barbero (les secteurs se chevauchent) mais il y a peut-être d'autres anomalies. Nous avons éliminé tous les doublons à moins de 100m les uns des autres. Il reste 35 anomalies de doublons éloignés l'un de l'autre. Il s'agit sans doute de reconnaissance erronée des numéros sur la photo de la Mappe par D. Barbero (exemple sur le secteur 44 une parcelle est lisible sur la carte comme la numéro 16875 mais elle est transcrit par sur par D.Barbero en 6871), l'autre explication possible une reconnaissance erronée des numéros sur les images Barbero par notre programme de retranscription.




Le décodage numéros de parcelles à partir des cartes de secteur Barbéro a abouti à reconnaître 11787 numéros de parcelles (et à les géoréférencer).

Les deux tables, celle du fichier excel du CD-Rom (enrichie par notre modélisation) et celle issue du géoréférencement des numéros de parcelles ont un index commun qui est le numéro de parcelle (dans quelques cas, une dizaine, une lettre est associée au numéro d'une parcelle, dans ce cas une décimale a été ajoutée au numéro de la parcelle par exemple la parcelle 17301 A a été numérotée 17301.02 cette décimale permet de ne pas confondre avec les décimales de 0.5 pour 1/2 et 0.33 pour 1/3 ). 

En fusionnant les deux fichiers, il est donc aisé de créer une seule carte vectorielle avec des données attributaires pour toutes les parcelles et une géolocalisation pour les 2/3 d'entre elles.

Mais il manque 6016 parcelles non géoréférencées, ce qui représente tout de même un tiers des parcelles.

Nous ne savons quels ont été les moyens et méthodes de Dominique Barbéro pour faire son travail de cartographie: Est-ce parce qu’il n’a pas pu lire certains numéros sur la Mappe originale ? Est-ce seulement l’édition des images de secteurs qui n’a pas affiché tous les numéros et toutes les parcelles ?

Toujours est-il que la proportion de parcelles non référencées pose un vrai problème pour l’objectif de cette étude qui est d’isoler les parcelles d’altitude.

\subsection{Geolocalisation précise sur carte Ign}

Nos trois couches :  photo de la Mappe (Raster), contour parcelle (Raster) et Parcelles Barbéro (vecteur), se superposent parfaitement.

Nous avons alors géoréférencé précisément la Mappe,  sur un fond de carte IGN en utilisant une variété de points de calage : \\
- les contours de la commune (en particulier les angles )\\
- les églises, chapelles, village...\\
- les cours d'eau\\
puis nous avons réutilisé ces points de calage pour géoréférencer les contours de parcelles et surtout la couche des parcelles.

\cartefigure
{figures/1 Mappe de Beaufort Georéférencée sur IGN.jpg}
{fig:MSsurIGN}
{Mappe de Beaufort Georéférencée sur IGN}

Le résultat est tout à fait satisfaisant, dans le sens où il n'y a pas de déformation exagérée de la Mappe de Beaufort quand on la transforme pour qu'elle se cale sur la carte moderne IGN

\subsection{Geolocalisation des parcelles manquantes}

Il reste, à ce stade 6016 parcelles non géoréférencées, mais nous connaissons les mas où elles se situent.

Les mas sont des territoires à l’intérieur du territoire de Beaufort, pour la plupart il s’agit d’un hameau ou groupement de maisons et les parcelles qui en dépendent, autrement ce sont des lieux dits. Toujours est-il que toutes les parcelles (le 17803) sont associées avec un « mas » et un seul.

Le CD-Rom Barbéro contenait une carte des mas, mais nous avons préféré reconstituer nous-mêmes les mas avec l'outil cartographique QGIS.

Nous avons donc créé une couche vectorielle des Mas, en créant des polygones qui entourent toutes les parcelles déjà géoréférencées.

Ce faisant nous avons constaté un certain nombre d'anomalies, en particulier des ensembles disjoints d'un même mas.

Nous avons alors créé un attribut "Mas2" qui est égale à "Mas" sauf dans certain cas, par exemple nous avons crée "Acrays 2" , "Ban 2", "Curtillet ?", "Dommellin ?", "Isle et Lacura ?", "Pecherie", "Plan ou Plains ?", "Praz" et "Territoire transféré à LaChapelle". Le principe était de n'avoir que pour chaque mas (Mas2) un territoire homogène.

\cartefigure
{figures/1 Carte des Mas.jpg}
{fig:Mas2}
{Carte des Mas reconstitués}

xxxxxxxxxxxxxxxx cas des parcelles hors de leur mas xxxxxxxxxxxx



\subsection{Ajout de l'altitude des parcelles}

Les cartes IGN fournissent aussi des couches de courbes de niveau, permettant de déterminer l'altitude de chaque point. Nous avons donc ajouté deux informations à la table attributaire des parcelles : l' altitude et une tranche d'altitude : \texttt{ [inf à 800m], [800-1000m], [1000-1200m], [1200-1400m], [1400-1600m], [1600-1800m] et [sup à 1800m]}.

Nous avons aussi ajouté une couche matérialisant les zones où l'altitude est supérieure à 1600m.

Il n’y a pas de définition précise des zones d’alpages par rapport aux zones d’hivernage permanent ou aux zones de pré-alpages. Dans le système agropastoral dit de « grande montagne » \footcite{ArbosMouthon2017} \footcite{Arbos1922}, il s’agit des pâturages utilisés entre la Saint XXX et le 24  juin (fête de la Saint Jean-Baptiste) et le 12 septembre (fête de la Saint XXX). 

Nous savons qu’à la fin de l’été, les grands troupeaux sont démantelés et dispersés dans des chalets capables de loger entre 5 et 10 vaches. Les vaches des grands troupeaux d’alpage pouvant appartenir soit au propriétaire (ou tenancier) de l’alpage soit au propriétaire (« client » de l’alpagiste) du chalet d’hivernage.

Suivant la richesse (en terre) des paysans, ils pouvaient posséder 3, 4 … 6 chalets intermédiaires entre la zone estivale et la zone d’hivernage finale en bas de vallée.

L’architecture des chalets varie en fonction de leur utilisation. Il est facile en se promenant dans la vallée du Beaufortain d’identifier à quelle altitude, les chalets d’hivernage disparaissent et sont remplacés par des granges, sans étable ni fenêtre ni espace d’habitation permanente. Par ailleurs une partie non négligeable des vieux chalets datent du 18e siècle. 

Cette approche empirique nous a conduits à fixer à 1600m l’altitude au-delà de laquelle on peut parler d’alpage. Nous avons identifié les derniers chalets d’habitat permanent à 1500m mais nous avons supposé que les prés ou pâturages avoisinant ces chalets étaient utilisés pour récolter le foin pour passer l’hiver et non pour la grande pâture.

Cette hypothèse est confirmée par l’exploitation actuelle des alpages, bien qu’il faudrait aussi prendre en compte l’histoire du climat. Le XVIIIe siècle était nettement plus froid qu’au XXe siècle où s’est mis en place l’agropastoralisme sous sa forme actuelle. Par conséquent la marge de 100m par rapport aux dernières habitations permanentes nous paraît aussi légitime de ce point de vue.




\subsection{Point d'étape : toutes les parcelles sont situées géographiquement}

Jusqu'à présent, le travail fut très technique, long et parfois fastidieux, mais cela nous a permis d'obtenir enfin une base de travail utilisable en Histoire.

Remarque : ce travail est déjà fait et disponible en ligne pour les communes de Hautes Savoie par les Archives Départementales de Haute-Savoie. Pour la Savoie, il est fait partiellement (CD Barbéro) et sur un nombre de communes limité. Sur la vallée du Beaufortain, seule la commune de Beaufort a été traitée au format Barbéro, rien n'est disponible pour les communes de Queige, Villard-sur-Doron et Hauteluce en dehors des documents d'origine de la Mappe. 


A ce stade nous avons donc une nouvelle table reprenant toutes les informations de la table initiale du CD Barbéro avec en plus pour chaque parcelle :\\
- son géoréférencement \\
- son altitude et sa classe d'altitude\\
- son occupation du sol plus utilisable que l'attribut nature du sol\\
- sa surface en ha (déduite de celle en m²)\\
- son mas (mas2) retravaillé pour n'avoir que des territoires d'un seul bloc \\
- des informations plus techniques comme le nom du secteur Barbéro sur lequel se situe la parcelle et ses coordonnées en pixel, et pour les parcelles n'en faisant pas partie la méthode qui a permis de les placer par rapport aux autres)



\section{Constitution d'une base de données enrichie}

Afin de rendre possible l'exploitation historique du cadastre, les données initiales ont été complétées par plusieurs informations nouvelles.

Les graphies manifestement variantes d'un même mas ont été harmonisées afin de constituer des ensembles territoriaux cohérents. Un attribut \texttt{Mas2} a ainsi été créé pour faciliter les traitements statistiques et cartographiques.

Les nombreuses dénominations cadastrales de la rubrique « Nature » ont également été regroupées dans une nomenclature simplifiée d'occupation du sol. Cette opération avait pour objectif de conserver la richesse descriptive des registres tout en permettant une lecture synthétique de l'organisation du territoire.

Enfin, chaque parcelle a été associée à une altitude obtenue par croisement avec un modèle numérique de terrain. Des classes altitudinales ont ensuite été définies afin de permettre l'étude de la répartition des propriétés, des usages du sol et des formes d'exploitation selon l'altitude.

\section{Résultat obtenu}

À l'issue de ce travail, chaque parcelle dispose :

\begin{itemize}
	\item des informations cadastrales issues des registres sardes ;
	\item d'une localisation géographique ;
	\item d'une appartenance à un mas ;
	\item d'une catégorie homogène d'occupation du sol ;
	\item d'une altitude et d'une classe altitudinale ;
	\item d'une superficie exprimée en hectares.
\end{itemize}

Cette base de données constitue le corpus utilisé dans l'ensemble des analyses présentées dans les chapitres suivants.

Les traitements techniques ayant permis la constitution de ce corpus sont présentés en détail dans l'annexe méthodologique.

\section{Méthodologie}

\subsection{Analyse de la table du CD-Rom : numérotation des parcelles}

Le CD-Rom Barbéro contient un ficher Excel : Beaufort-FICHIER-ATLAS-OK.xls contenant 17803 parcelles  (avec en plus, une ligne d’en-têtes des colonnes) :

•	Numéro : il s’agit des numéros de parcelles a priori numériques dont le plus grand numéro est 17691. Il devrait a priori avoir 17691 parcelles sur la commune de Beaufort du Cadastre Sarde mais on constate immédiatement que le nombre de lignes ne correspond pas au nombre de parcelles

•	CF : Chef-lieu « O » ou « null » ; signification à investiguer ultérieurement

•	E : Ancien fond ecclésiastique « O » ou « null » ; signification à investiguer ultérieurement

•	I : « O » ou « null » ; signification à investiguer ultérieurement (dans la notice livrée avec le CD-Rom il est fait mention à cet emplacement d’une colonne nommée « N » (pour Noble)

•	Propriétaires : nous avons là des noms et prénoms des propriétaires avec parfois des compléments d’information par exemple « MONOD Jeanne soit Popelin veuve d'Antoine GACHET » ou « TARIN Nicolas et Humbert son frère» et aussi des noms de communautés par exemple « La communauté pour l'usage commun » ou « Chapelle de Roselend »

•	Statut : cette colonne contient du texte avec des valeurs répété : « Bourgeois » « communier » … il s’agit d’une information sur le statut du propriétaire a priori tel qu’il est inscrit sur le registre associé au cadastre

•	IND : parcelle en indivision, elle contient des valeurs « null » et quelque « O »

•	Mas : 

•	Nature : cette colonne contient du texte avec des valeurs récurrentes : « Roc », « Pâturage » ; « Marais et pâturage » … qui sont des descriptions de la nature du sol

•	DB : cette colonne contient des valeurs numériques 0, 1 , 2 ou 3 qui correspond aux quatre « degrés de bonté ». Le degré de bonté est une évaluation de la qualité (ou du rendement) du sol, faite par les agents du duché de Savoie en charge du relevé du cadastre, il a un but fiscal

•	m² : il s’agit sans doute de la surface de la parcelle en m² ; le registre du cadastre contenant de valeur de superficie selon une double mesure : la mesure de Savoie (Journaux, toises, pieds) et mesure de Piémont (Journaux, Tables, Pieds). 1 journal de Savoie équivaut à 0,3685 ha, celui du Piémont lui est légèrement supérieur. Cette conversion a été réalisée sans doute par le créateur du fichier D.BARBERO. (remarque : l’ensemble de la commune actuelle couvre 149,53 klm²  mais la partie transférée à l’époque de l’établissement de la Mappe Sarde est revenue à la commune actuelle de Beaufort sur Doron, la somme des surfaces du fichier D.BARBERO est égale à 139.648.325 m² soit 139,6 klm² ce qui est cohérent, on peut donc valider que la surface sur fichier Excel est en m² d'autant plus qu'à l'époque de l'établissement de la Mappe une partie des parcelles a été rayée car réattribuée à la paroisse de LaChapelle, elles ont été depuis réintégrées à la commune de Beaufort, dont il est normal que la superficie calculée par Barbéro soit inférieure à la superficie de la commune actuelle de Beaufort-sur-Doron) 

•	Avoine : des valeurs null ou numérique s’étalant de 13,5 à 22

•	Seigle : des valeurs null ou numérique s’étalant de 0,5 à 4 (les valeurs 0,5 sont stockées au format texte « 1 / 2 » nous les avons changées en 0,5

•	F. Bœuf : sous bœuf est regroupé le nombre de bovins, la colonne contient des valeurs numériques de 0,5 à 9 puis pour les valeurs supérieures le nombre est suivi de la lettre « L » : « 10 L », « 12 L », « 100 L » …

•	F. Cheval : sous Cheval est regroupé le nombre d’équidés de toute sorte (âne, mule, mulet cheval ) . La colonne contient des valeurs numériques de 0,5 à 9 puis pour les valeurs supérieures le nombre est suivi de la lettre « L » : « 10 L », « 12 L », « 100 L » …

•	Fascines : (fagot de petit bois ?) valeur numérique non investiguée pour l’instant

•	Chevron : valeur numérique sans doute bois d'œuvre pour charpente

•	Latte : valeur numérique sans doute bois d'œuvre pour des planches

•	Pièce 1 : valeur numérique non investiguée pour l’instant

•	Pièce 2 : valeur numérique non investiguée pour l’instant

•	N° Grief : champs pouvant contenir une ou plusieurs valeurs numériques, il s’agit de numéro de référence des contestations faites auprès des autorités relativement à la parcelle cadastrée. En principe il faudrait normaliser ce champ et créer une logique de relation 1-n entre une parcelle et les griefs associés, cependant il n’y a que 22 cas où une parcelle est associée à 2 ou 3 griefs, et l’étude des griefs est très secondaire dans notre étude.

•	Grief : champ texte contenant une typologie selon les types de griefs

•	Problèmes : descriptif sommaire du grief


Degré de Bonté :

Dans l’attribution des degrés de bonté, chaque type d’occupation du sol comporte les trois classes ; la qualité de la parcelle pouvant être bonne (1), moyenne (2) ou mauvaise (3). Par exemple pour la paroisse de Landry en vallée de Tarentaise (Savoie) la « note des estimes en degré avec leurs catégories… » nous informe sur la production qu’un terrain doit fournir pour se voir attribuer tel ou tel degré de bonté. Un pâturage doit produire annuellement deux quintaux de foin pour obtenir le degré de bonté 1, un quintal et demi pour le 2 et un quintal pour le 3. Pour les champs, la mesure est en bichet et il en faut vingt-neuf répartis en froment, seigle et farine pour répondre aux critères du degré de bonté 1, vingt-six pour le 2 et vingt-trois pour le 3, et ainsi de suite pour chaque usage du sol. Pour le bâti, les estimateurs attribuent le degré de bonté d’une terre voisine, de préférence appartenant au même propriétaire. Cette méthode d’évaluation des terrains fixe des niveaux d’exigence précis basés sur le rendement de la parcelle. Ainsi, les parcelles classées en 1 sont de bonnes terres productives, car les critères exigés pour avoir ce degré de bonté sont du même niveau. Le fait que chaque nature de sol comporte l’intégralité des classes (1, 2 et 3) et que les critères d’évaluation répondent à un même niveau d’exigence permet une comparaison entre deux terrains du même degré de bonté. \footcite{Baud2010} 


Mesure de Savoie :
« La base de ce système est le « pied de chambre de 0,033m qui élevé au carré donne le « pied de cadastre de 0.92m². un carré de 8 pieds forme la toise carrée (7.370m²) et 400 toises carrées équivalent à un journal de Savoie dont la valeur exacte est de 29.4838 ares. On compte pour les approximations 3 journaux à l’hectare » 	\footcite{GUICHONNET1955}


Mesure de Piemont :

« Les cadastreurs piémontais se servirent des unités de leur pays et qui dérivaient du pied liprand de 0,513 m. Il fallait 6 pieds pour faire un trabuc (3,082 m) et 2 trabucs pour une perche (6,164 m). Ces mesures linéaires servaient de base aux unités agraires suivantes : le pied de table (rectangle de 1 x12 pieds liprands de côté) ou 3,167 m². La table, ou perche carrée vaut 12 pieds de table (38,009 m) et 100 tables, ou 400 trabucs carrés, donnent un journal de Piémont. La mesure de Piémont est à préférer lorsqu’on utilise l’ancien cadastre. Pour les usagers soucieux de l’exactitude des contenances, elle est évidemment meilleure car elle résulte des opérations sur le terrain ; pour l’historien, elle est commode car la conversion des tables en journaux repose sur la base 100. » Idem


Mas :

«On a la preuve, par les reconnais; féodales, que les mas sont antérieurs de trois ou quatre siècles au cadastre, et que les « trabucants » piémontais n’ont fait qu’enregistrer un état de fait profondément et obscurément enraciné dans la tradition paysanne. Les mas appartiennent à plusieurs propriétaires, nobles ou roturiers, et le régime foncier de 1730 recouvre une réalité agraire très antérieure; ils chevauchent parfois plusieurs communes et ne portent pas toujours le nom des villages ou des lieux-dits qu’ils renferment, mais un toponyme d’étymologie romaine ou germanique. Leur étendue comprend rarement une seule zone de végétation, mais englobe différents types de terroirs et semble bien représenter une unité élémentaire d'exploitation, domaine d’un seul individu. » Idem p296

a) Nettoyage des données

•	Pb colonne « Numero »

En fin de fichier il y a 13 lignes avec des valeurs alphanumériques : 12 parcelles ont une lettre à la suite du numéro de parcelle et 1 parcelle a une lettre sans numéro :

17160 B;;;;BOUCHAGE Jean Baptiste;..\\
17161 A;;;;BUGAND Claude;… \\
17301 A;;;;MARTIN Claude;… \\
17302 B;;;;MARTIN Claude;… \\
17303 C;;;;La communauté pour l'usage commun;Communaux;… \\
17304 D;;;;La communauté pour l'usage Commun ;Communaux ;… \\
17305 E;;;;La communauté pour l'usage commun;Communaux;…\\
17306 F;;;;La communauté pour l'usage commun;Communaux;… \\
17307 G;;;;La communauté pour l'usage commun;… \\
17308 H;;;;La communauté pour l'usage commun;…\\
17309 I;;;;La communauté pour l'usage commun;\\
17310 K;;;;La communauté pour l'usage commun;\\
A;;;;TOCCAT Daniel;Communier;;;;3;1543;;;;;;;;;;;;\\

Nous avons créé une nouvelle colonne : Complément_numéro.

Les lettres des 13 lignes sont déplacées dans cette nouvelle colonne, et la ligne sans numéro se voit attribuer un numéro « 0 ».

Il y a aussi des parcelles dont le numéro est suivi de 1/2 , 1/3 ou 1/4 , ces valeurs ont été remplacées par des décimales (0.5, 0.3, 0.25) ajoutées au numéro de parcelle.

Il reste à comprendre pourquoi le nombre de lignes est significativement différent de la valeur maximum de numéro de parcelle. Mais nous allons d’abord importer ce fichier dans une base de données de travail.

•	Pb colonnes « CF » , « E », « I » , « IND »

Les valeurs « O » dans les colonnes « CF » , « E », « I » , « IND » ont été remplacées par des « Oui » afin de permettre une caractérisation booléenne.

•	Pb colonne « m² »

147 ;;;;VIALLET Jacques fils de Benoit;Communier;;Praz;Chenevier;1;\#VALUE!;21;3;;;;;;;;;;\\
La valeur \#value est remplacé par la valeur « null » \\
•	Pb colonne « Nature »

La colonne Nature est longue et contient beaucoup de double description, par exemple « bois et pâturage » ou « prés et teppes », une option de normalisation aurait été de décomposer en une relation 1-N : une parcelle pouvant être reliée alors à deux natures « bois et « pâturage » ; mais n’étant pas certain de comprendre précisément l’intention du géomètre-cadastreur, nous avons pris l’option de laisser pour l’instant le champ nature tel quel, mais en revanche de créer un concept « Occupation_Sol » pour regrouper les nombreuses « Natures » en ensemble homogène en prenant pour principe lorsqu’une Nature est « multiple » de prendre la première comme la Nature principale ; par exemple « « bois et pâturage » » aura la valeur « Bois » pour « Occupation_Sol »

•	Pb colonne « Mas »

Pour les Mas il y a des mas dont les libellés se ressemblent, par exemple : « Coutafalliat » et » Coutassalliat » (les » s » et les «f » dans la graphie du XVIIIe), mais comme les numéros de parcelles ne sont pas forcément continus entre les deux orthographes, nous avons décidé de laisser les libellés de mas tel quel en attendant un travail futur de représentation cartographique des mas et parcelles.

•	Pb colonnes « Avoine », «Seigle », « Boeuf », « Cheval », « Fascines » , « Chevron » , « Latte» :

Les champs « Avoine », «Seigle », « Boeuf », « Cheval », « Fascines » , « Chevron » , « Latte» décrivent soit une quantité de production d’Avoine de seigle, de fagots de bois (Fascines ?) et c’est logique qu’elle soit associée à une certaine parcelle précise, mais pour les autres il s’agit d’un nombre de bêtes ; cette dernière valeur pose question, car cela devrait:

•	Être associé au propriétaire et correspondre à la taille de son troupeau (cela poserait alors d’autres questions, car pour la gestion des alpages (en été) les troupeaux sont constitués de bêtes appartenant en propre et des bêtes en location) ; cela ne semble pas être le cas, car pour un propriétaire possédant plusieurs parcelles le nombre de bêtes varie selon une logique non claire ; exemple :
Parcelles	Propriétaires	Mas	Nature	DB	Surface	Boeuf\\
7806	BLANC Jean Baptiste	Val de Treicol	Grange	2	67	180\\
7807	BLANC Jean Baptiste	Val de Treicol	Grange	2	181	180\\
7808	BLANC Jean Baptiste	Val de Treicol	Pâturage	2	31130	180\\
7774	BLANC Jean Baptiste	Val de Treicol	Grange	2	206	160\\
7776	BLANC Jean Baptiste	Val de Treicol	Pâturage et rocher	2	86210	160\\
8331	BLANC Jean Baptiste	Couecle	Pâturage et broussailles	2	65791	150\\
8323	BLANC Jean Baptiste	Couecle	Pâturage	2	147294	120\\
8324	BLANC Jean Baptiste	Couecle	Grange	2	79	120\\
7802	BLANC Jean Baptiste	Val de Treicol	Pâturage	3	10135	50\\
7803	BLANC Jean Baptiste	Val de Treicol	Pâturage	3	22085	30\\

•	Ou être associé à une ou plusieurs parcelles : il est évident qu’un troupeau qui « remue » change de parcelle régulièrement.
Nous avons donc conservé ces champs tels quels comme attribut de la parcelle en attendant de comprendre comment cette donnée fonctionne (il faudra sans aucun doute revenir aux registres originaux manuscrits associés au Cadastre Sarde conservé aux Archives Départementales de Savoie)

•	Pb colonnes « Avoine » et « Seigle »

Par exemple dans la ligne :

5299 ;;;;DUITTOZ GENIS;Communier;;Bersend;Champ;3;377;14 1/2; 1/2;;;;;;;;;;\\
Il est inscrit 14 1/2 pour 14,5 et 1/2 pour 0.5

En ouvrant le fichier avec Excel, les valeurs sont bien comprises par le logiciel comme des nombres décimaux, sauf lorsque la valeur est 1/2 qui reste comprise comme du texte

Les 668 cas des colonnes « Avoine » et « Seigle » ont été corrigés.

•	Pb colonnes Bœuf et Cheval

Il existe des « L » qui suivent les valeurs supérieures à 10 ; je ne sais pas à quoi cela correspond (sans doute pas l’unité de mesure Livre. Ces « L » sont supprimés afin d’avoir des colonnes numériques. (Il faudra tout de même vérifier ce qui est inscrit dans le registre.
Il y a aussi un cas où la cellule Bœuf contient une virgule ; elle est remplacée par « Null»

Pour la colonne Cheval il y a aussi 4 cas de « 1/2» remplacés par « 0.5»

•	Pb colonnes No Grief, Grief, Problème

Il existe des lignes dont le champ « Problèmes» n’a pas de numéro de Grief associé sans numéro., nous avons ajouté une numérotation allant de 800 à 952 et la description « Grief sans numéro ».

•	divers :
Nous avons aussi modifié certains titres de colonnes (remplacement des é , ², œ, espace…) et ajouté une colonne » Analyse » (vide) qui servira pour stocker les résultats des analyses des numéros de parcelles.

Désormais nous avons un fichier csv contenant des champs propres et homogènes :

Analyse	Texte\\
Numero	Numerique\\
Complement_numero	Texte\\
CF	booléen (Oui/Non)\\
E	booléen (Oui/Non)\\
I	booléen (Oui/Non)\\
Proprietaires	Texte\\
Statut	Texte\\
IND	booléen (Oui/Non)\\
Mas	Texte\\
Nature	Texte\\
DB	Numerique\\
Surface	Numerique\\
Avoine	Numerique\\
Seigle	Numerique\\
Bœuf	Numerique\\
Cheval	Numerique\\
Fascines	Numerique\\
Chevron	Numerique\\
Latte	Numerique\\
Pièce1	Numerique\\
Pièce2	Numerique\\
No_Grief	Texte\\
Grief	Texte\\
Problemes	Texte\\

\subsection{Constitution d’une table attributaire enrichie}

Un travail de modélisation de la table contenue dans le CD Barbero est disponible en annexe: ce travail permet de distinguer les entités : Mas, Propriétaire et Nature

Au niveau des Mas nous avons par exemple fusionné le mas "Coutassallat" dans le mas "Coutafaillat" ainsi que le mas "Treverse" dans le mas "Traverse"

Au niveau des nombreuses "natures", nous avons créé une entité "Occupation du Sol" afin de regrouper les 440 natures différentes.

Regroupement des natures cadastrales

Les registres de la Mappe sarde décrivent chaque parcelle au moyen d'une rubrique dite « nature ». Cette information est particulièrement riche et reflète la diversité des milieux exploités par les communautés montagnardes du Beaufortain. Outre les catégories classiques telles que les champs, prés ou bois, les estimateurs distinguent de nombreuses situations intermédiaires : pâturage sur roc, roc et pâturage, bois et broussailles, rocher broussailles, glières, rippes, marais, rivages, etc. Au total, plusieurs dizaines de libellés différents apparaissent dans les registres.

Cette richesse descriptive constitue un atout pour comprendre la perception et l'évaluation des terres par l'administration sarde. En revanche, elle rend difficile toute analyse statistique ou cartographique globale. La multiplication des catégories produit des tableaux peu lisibles et empêche de dégager les grandes structures de l'occupation du sol à l'échelle de la communauté.

Afin de permettre les traitements statistiques, les différentes natures ont donc été regroupées en un nombre limité de catégories fonctionnelles. Ce regroupement ne vise pas à remplacer les dénominations originales, qui demeurent conservées dans la base de données, mais à faciliter l'analyse des grandes composantes du territoire.

La méthode retenue repose sur une hiérarchie fondée sur l'usage économique principal de la parcelle. Lorsqu'un libellé comporte le terme « pâturage », la parcelle est classée dans la catégorie « Pâturage », quelle que soit la présence éventuelle d'autres éléments (roc, bois ou broussailles). Ce choix s'explique par l'importance économique de l'activité pastorale dans le Beaufortain et par le fait que de nombreuses parcelles d'alpage associent naturellement herbe, affleurements rocheux et végétation ligneuse. (remarque : nous n'avons pas estimé fiable l'idée que l'estimateur sarde ait marqué une prédominance avec l'ordre des môts dans son descriptif, ainsi "pâturage et bois" et sera pas différentié de "bois et pâturage").

Les parcelles ne comportant pas de pâturage sont ensuite examinées pour identifier les terres destinées aux cultures (champs, jardins et autres terres cultivées), puis les prés de fauche. Viennent ensuite les espaces boisés, puis les broussailles et bruyères. Enfin, les terrains essentiellement minéraux (roc, roche, rocher, ravines, murgers, glières, etc.) sont regroupés dans la catégorie « Roc ». Une catégorie spécifique est également conservée pour le bâti ainsi que pour les surfaces en eau, marais et rivages. Cette logique hiérarchique est biaisé et orientée avec le parti pris de a) l'importance économique majeur des pâturages pour le Beaufortain en 1730, puis b) les champs permettant des cultures vivrières et c) les prés de fauche de fourrage pour l'hiver, nous avons placé d) les bois avant e) les broussailles (ces dernières ne sont pas du tout improductives car elles permettent de nourrir des caprins et même de fournir une alimentation en hiver pour tout le bétail ) enfin nous avons f) les roc g) les marais... h) le bâti

Cette classification présente l'avantage de restituer les grandes fonctions économiques du territoire tout en restant fidèle à la logique des registres sardes. Elle permet notamment de distinguer les espaces agricoles, les ressources fourragères destinées à l'hivernage du bétail, les pâturages d'altitude, les ressources forestières et les terrains à faible potentiel productif. L'analyse croisée de ces catégories avec l'altitude et le degré de bonté permet ainsi la mise en évidence de l'organisation agropastorale de la communauté de Beaufort au XVIIIe siècle.










